\subsection{有理数有序群}

我们在有理数集 Q 上定义过序关系 \(x \leqslant y\) ；我们已经看到，这个序使 Q 成为一个全序集，并且它与 Q 的加法群结构相容，即对每个 \(z \in Q\) ，关系 \(x \leqslant y\) 等价于 \(x + z \leqslant y + z\) （也就是说，该序在平移下不变）。我们回忆一下记号（它在任何全序群中都使用）

\[
\begin{array}{l} x ^ {+} = \sup (x, 0), \\ x ^ {-} = \sup (- x, 0) = (- x) ^ {+}, \\ | x | = \sup (x, - x); \end{array}
\]

\(|x|\) 称为 \(x\) 的绝对值，我们有

\[
x = x ^ {+} - x ^ {-}, \quad | x | = x ^ {+} + x ^ {-}
\]

以及三角不等式

\[
| x + y | \leqslant | x | + | y |,\tag{1}
\]

以及不等式

\[
\mid | x \mid - \mid y \mid \mid \leqslant \mid x - y \mid\tag{2}
\]

它是 (1) 的直接推论；此外还有

\[
\left| x ^ {+} - y ^ {+} \right| \leqslant | x - y |.\tag{3}
\]

关系 \(x \geqslant 0\) 、\(x = x^{+}\) 、\(x^{-} = 0\) 、\(|x| = x\) （相应地 \(x \leqslant 0\) 、\(x = -x^{-}\) 、\(x^{+} = 0\) 、\(|x| = -x\) ）彼此等价。关系 \(|x| = 0\) 等价于 \(x = 0\) ；若 \(a \geqslant 0\) ，则关系 \(|x| \leqslant a\) 等价于 \(-a \leqslant x \leqslant a\) ，而关系 \(|x| \geqslant a\) 等价于“ \(x \geqslant a\) 或 \(x \leqslant -a\) ”。对所有 \(x, y\) ，在 \(\mathbf{Q}\) 中我们有

\begin{enumerate}
\def\labelenumi{(\arabic{enumi})}
\setcounter{enumi}{3}
\tightlist
\item
\end{enumerate}

\[
\sup (x, y) + z = \sup (x + z, y + z),\tag{4}
\]

\[
\inf (x, y) = - \sup (- x, - y),\tag{5}
\]

而作为特殊情形，

\begin{enumerate}
\def\labelenumi{(\arabic{enumi})}
\setcounter{enumi}{5}
\tightlist
\item
\end{enumerate}

\[
\sup (x, y) = x + (y - x) ^ {+} = x + (x - y) ^ {-},\tag{6}
\]

\[
\inf (x, y) = x - (y - x) ^ {-} = x - (x - y) ^ {+},\tag{7}
\]

最后，令 \(Q_{+}\) 表示 \(\geqslant0\) 的有理数组成的集；于是我们有

\begin{enumerate}
\def\labelenumi{(\arabic{enumi})}
\setcounter{enumi}{7}
\tightlist
\item
\end{enumerate}

\[
\mathbf {Q} _ {+} + \mathbf {Q} _ {+} \subset \mathbf {Q} _ {+},\tag{8}
\]

\[
\mathbf {Q} _ {+} \cap (- \mathbf {Q} _ {+}) = \{\mathrm{o} \},\tag{9}
\]

\[
\mathbf {Q} _ {+} \cup (- \mathbf {Q} _ {+}) = \mathbf {Q},\tag{10}
\]

关系 \(x \leqslant y\) 等价于 \(y - x \in \mathbf{Q}_{+}\) 。

我们将利用这个序来定义 Q 上一个与其加法群结构相容的拓扑。

